\documentclass[11pt]{article}
\usepackage[margin=1in]{geometry}
\usepackage{amsmath,amssymb,amsthm,mathtools}
\usepackage{graphicx,microtype,xcolor}
\usepackage[colorlinks=true,linkcolor=blue!55!black,citecolor=blue!55!black,urlcolor=blue!55!black]{hyperref}

\newtheorem{theorem}{Theorem}
\newtheorem{lemma}{Lemma}
\newtheorem{remark}{Remark}
\newcommand{\T}{\mathbb T}
\newcommand{\Hd}{\mathcal H}

\title{Hausdorff Isolation and a Stopping-Scale Barrier\\
for Rajchman Multiplier Norms}
\author{Automated mathematics run \texttt{fa\_banach\_001}, lane 7}
\date{17 August 2026}

\begin{document}
\maketitle

\begin{center}
\fbox{\parbox{0.92\textwidth}{\textbf{Status.} Substantial partial result,
likely valid, pending human review.  The isolation lemma asked about in arXiv
v1 extends completely from Minkowski to Hausdorff dimension.  For the stronger
question in the live arXiv v2, the mountain--river estimate extends to
variable-depth Hausdorff stopping covers.  The unresolved step is to
synchronize their good vertices at one physical scale.}}
\end{center}

\section{Version-sensitive source question}

Ohrysko, Wojciechowski, and Zawalski's cached arXiv v1 asks whether their
measure-isolation lemma can use Hausdorff rather than Minkowski dimension
\cite{OWZv1}.  The live arXiv v2, posted 9 August 2026 under a revised title,
substantially strengthens that question: it asks whether lower Minkowski
dimension can be replaced by Hausdorff dimension in the full multiplier-norm
theorem \cite[Open Problem 4]{OWZv2}:

\begin{figure}[ht]
\centering
\includegraphics[width=0.94\textwidth]{figures/open_problem_crop.png}
\caption{The current Hausdorff-dimension question in arXiv v2, PDF p.~29.}
\end{figure}

For $0<\beta\le d$ and a finite positive Borel measure $\nu$ on $\T^d$, put
\[
 h_\beta(\nu)=\sup\{\nu(E): E\subset\T^d\text{ Borel and }
                         \Hd^\beta(E)=0\}.
\]
Using $\Hd^\beta(E)=0$ includes every $E$ with
$\dim_H E<\beta$ and gives the natural endpoint formulation.

\section{Hausdorff isolation lemma}

For a set $U\subset\T^d$, write
$U^{(\delta)}=\{x:\operatorname{dist}_\infty(x,U)<\delta\}$.

\begin{theorem}[Hausdorff isolation]\label{thm:isolation}
Let $\nu,\eta$ be mutually singular finite positive Borel measures on
$\T^d$.  For every $\tau>0$ and $\delta_0>0$ there are
$0<\delta<\delta_0$ and a finite family $\mathcal D$ of dyadic cubes such
that, with $U_{\mathcal D}=\bigcup_{Q\in\mathcal D}Q$,
\begin{align*}
 \ell(Q)&<\delta &&(Q\in\mathcal D),\\
 \sum_{Q\in\mathcal D}\ell(Q)^\beta&<1,\\
 \nu(U_{\mathcal D})&>h_\beta(\nu)-\tau,\\
 \eta(U_{\mathcal D}^{(\delta)})&<\tau.
\end{align*}
In particular, the v1 isolation question has an affirmative answer.
\end{theorem}

\begin{proof}
Choose a Borel $E$ with $\Hd^\beta(E)=0$ and
$\nu(E)>h_\beta(\nu)-\tau/3$.  Mutual singularity supplies a Borel set $S$
with $\nu(\T^d\setminus S)=0$ and $\eta(S)=0$.  By inner regularity choose
a compact
\[
 K\subset E\cap S,\qquad \nu(K)>h_\beta(\nu)-2\tau/3.
\]
Then $\Hd^\beta(K)=0$ and $\eta(K)=0$.  Outer regularity gives an open
$O\supset K$ with $\eta(O)<\tau$.

We use the elementary equivalence of ordinary and dyadic Hausdorff content.
Indeed, a cube of side $r$ can be covered by at most $C_d$ dyadic cubes of
side at most $2r$.  Since $\Hd^\beta(K)=0$, for every $\varepsilon>0$ and
every scale bound there is therefore a finite dyadic cover $\mathcal D$ of
$K$, as fine as desired, with
\[
                 \sum_{Q\in\mathcal D}\ell(Q)^\beta<\varepsilon.
\]
Take $\varepsilon<1$.  Compactness of $K\subset O$ gives
$r=\operatorname{dist}_\infty(K,\T^d\setminus O)>0$ (with the evident
interpretation if $O=\T^d$).  Choose the cover so fine that every cube lies
in $K^{(r/4)}$ and has side less than
$\min(\delta_0/2,r/8)$.  Now select $\delta$ strictly larger than all side
lengths and smaller than $\min(\delta_0,r/2)$.  After making the preliminary
cover finer if necessary, its $\delta$-enlargement lies in $O$.  Hence the
last displayed $\eta$-bound holds, while the cover of $K$ gives the required
$\nu$-bound.  The content and scale bounds hold by construction.
\end{proof}

The proof uses only Hausdorff content and regularity.  Thus the v2 obstruction
does not occur in separating a thin one-signed component from the mutually
singular remainder.

\section{Mountain--river combinatorics on a Hausdorff stopping cover}

The next statement shows that the other combinatorial heart of the source
also survives variable scales.  Consider the rooted $2^d$-ary dyadic tree.
For a positive measure let $m(Q)$ be the mass below a vertex $Q$.  Fix
$-d<\alpha<0$ and call $Q$ \emph{turbulent} if every child $Q'$ satisfies
$m(Q')<2^\alpha m(Q)$.

\begin{lemma}[Variable-depth terminal estimate]\label{lem:stopping}
Let $\mathcal L$ be a finite dyadic antichain whose union carries a positive
measure of total mass $M$, and suppose every $R\in\mathcal L$ has depth
$|R|\ge K$.  Fix $0<\rho<1$ and put $\gamma=-\alpha\rho$.  If
$\gamma>\beta$, then the total mass of leaves $R$ whose root-to-$R$ path
contains at least $\rho|R|$ turbulent vertices is at most
\begin{equation}\label{eq:bad}
 M\,2^{-(\gamma-\beta)K}
       \sum_{R\in\mathcal L}2^{-\beta|R|}.
\end{equation}
\end{lemma}

\begin{proof}
If $S(R)$ is the number of turbulent ancestors of $R$, monotonicity at all
other steps and the turbulent inequality give
\[
                    m(R)\le M,2^{\alpha S(R)}.
\]
For a leaf counted on the left of \eqref{eq:bad}, this is at most
$M2^{-\gamma|R|}$.  Consequently
\[
 \sum_{R\ \mathrm{bad}}m(R)
 \le M\sum_{R\ \mathrm{bad}}2^{-\gamma|R|}
 \le M2^{-(\gamma-\beta)K}
            \sum_{R\in\mathcal L}2^{-\beta|R|},
\]
as claimed.
\end{proof}

The numerical condition $\gamma>\beta$ is exactly
$\alpha\rho+\beta<0$, the condition in the source's fixed-generation
mountain--river lemma.  Since a compact $\Hd^\beta$-null set has arbitrarily
fine stopping covers with arbitrarily small
$\sum_R2^{-\beta|R|}$, Lemma~\ref{lem:stopping} says that almost all of its
mass travels through many non-turbulent vertices.  Thus both the isolation
and high-turbulence estimates admit Hausdorff versions.

\section{The remaining synchronization gap}

The v2 proof next selects one generation $n$ at which a positive fraction of
the mass is non-turbulent.  It then uses a single translated Hardy atom of
side $2^{-n}$, simultaneously testing every selected cube.  A Hausdorff
stopping cover provides terminal depths $|R|$ with no bounded ratio.  The
good ancestors supplied by Lemma~\ref{lem:stopping} can therefore be spread
over arbitrarily many generations.  A large measure-weighted total number of
good vertices does not imply that any one generation carries a fixed
fraction of $M$.

Summing one normalized Hardy atom for each active scale pays the $\ell^1$
sum of its atomic coefficients.  Square-function organization improves this
only to an $\ell^2$ aggregation, still losing when the captured mass is spread
over many scales.  Refining all stopping cubes to one generation is equally
insufficient: a cube of depth $j$ creates $2^{d(k-j)}$ descendants at depth
$k$, whereas Hausdorff content controls only $2^{-\beta j}$.

Eight focused proof/counterexample attempts were made.  Besides the two
proved statements above, the most concrete negative route is conditional:
if a singular Hausdorff-thin probability $\sigma$ induced a compact map
$T_\sigma:H^1\to L^1$, then generic small translations would give
$\sigma\perp\tau_h\sigma$ but
\[
 \|T_{\sigma-\tau_h\sigma}\|
   =\|(I-\tau_h)T_\sigma\|\longrightarrow0,
\]
contradicting any Hausdorff lower bound.  Existence of such a compact thin
convolution operator was not established.  Conversely, no multiscale atom
with scale-independent $H^1$ cost was found.  These are genuine remaining
tasks, so the current v2 question is not claimed solved.

\section{Novelty and verification boundary}

The exact question, arXiv id and titles, run indexes, and arXiv searches for
Hausdorff-dimension Rajchman multiplier results were checked on 17 August
2026.  The live v2 itself, posted eight days earlier, still lists the stronger
question as open; no later answer was located.  Theorem~\ref{thm:isolation}
is elementary but closes the precise v1 question.  Lemma~\ref{lem:stopping}
is the substantive partial advance for v2 and sharply identifies the missing
scale-synchronization step.  Human review should focus on the dyadic-content
conversion in Theorem~\ref{thm:isolation} and on whether a known stopping-time
$H^1$ construction can overcome the final gap.

\begin{thebibliography}{9}
\bibitem{OWZv1}
P. Ohrysko, M. Wojciechowski, and B. Zawalski,
\emph{On the Wiener--Pitt phenomenon for algebra of Rajchman multipliers on
Hardy space}, arXiv:2206.06958v1, 2022.

\bibitem{OWZv2}
P. Ohrysko, M. Wojciechowski, and B. Zawalski,
\emph{The $H^1(\T^d)$-multiplier norm of Rajchman measures, Zafran-type
phenomena, and Walsh--Stieltjes series}, arXiv:2206.06958v2, 2026.
\end{thebibliography}

\end{document}
